\chapter{Linear maps and change of basis}\label{ch:maps}

Chapter 2 called a matrix ``a recipe'' without defining the recipe kind. This chapter names it: a matrix is a \emph{linear map} written in coordinates — a function between vector spaces that respects the two operations. That single definition explains why the algebra of Chapter 2 exists at all (matrix multiplication is composition), why coordinates are a design decision (Chapter 3), and why two different matrices can be \emph{the same map wearing different labels}. After it you will be able to test a map for linearity, read off its matrix, move the same map between bases, and place the four classic planar maps — rotation, reflection, projection, shear — in a permanent mental gallery.

\section{Linear maps: definition and examples}\label{sec:linmaps}

A map $T:V\to W$ between vector spaces is \emph{linear}\index{linear map} if it respects addition and scaling:
\[
  T(\mathbf{u}+\mathbf{v}) = T(\mathbf{u})+T(\mathbf{v}),\qquad
  T(c\mathbf{v}) = c\,T(\mathbf{v}).
\]
Equivalently one line: $T(a\mathbf{u}+b\mathbf{v})=aT(\mathbf{u})+bT(\mathbf{v})$. The consequences matter more than the check: a linear map is completely determined by what it does to a basis (everything else is forced by the two rules), and it must send $\mathbf 0$ to $\mathbf 0$ (put $c=0$). That second clause is a free disqualifier — any map that moves the origin is out before the real test starts.

The family is bigger than it looks: matrix-by-vector $T(\mathbf{x})=A\mathbf{x}$ is linear (that is the distributive law); differentiation of polynomials is linear ($\frac{d}{dt}(p+q)=p'+q'$); the shift $(x_1,x_2,\dots)\mapsto(0,x_1,\dots)$ on sequences is linear; rotation, reflection, projection, shearing of the plane are all linear — Section~\ref{sec:gallery} computes their matrices. Non-examples: translation $\mathbf{x}\mapsto\mathbf{x}+\mathbf{b}$ moves $\mathbf 0$; $\mathbf{x}\mapsto\|\mathbf{x}\|$ fails scaling with $c=-1$; $\mathbf{x}\mapsto x_1^2$ fails everything. Figure~\ref{fig:affine} keeps the three cases apart in one look.

\begin{figure}[htbp]
\centering
\begin{tikzpicture}[scale=0.9, >=stealth]
  % linear: line through origin stays through
  \begin{scope}
    \draw[->,gray] (-1.5,0)--(1.5,0); \draw[->,gray] (0,-1.2)--(0,1.4);
    \draw[accent,very thick] (-1.2,-0.7)--(1.2,0.9);
    \node at (0,-1.7) {\small linear $T$: lines through $\mathbf 0$ stay};
    \node[align=center] at (0,1.9) {\small $\mathbf 0\mapsto\mathbf 0$,\quad grid squares\\ \small stay parallelograms};
  \end{scope}
  % affine: shifted
  \begin{scope}[xshift=5cm]
    \draw[->,gray] (-1.5,0)--(1.5,0); \draw[->,gray] (0,-1.2)--(0,1.4);
    \draw[amber,very thick] (-1.2,-1.0)--(1.2,0.6);
    \fill (0,0) circle (1.2pt) node[below left,ink]{\tiny $\mathbf 0\notin$ img};
    \node at (0,-1.7) {\small affine $\mathbf x\mapsto A\mathbf x+\mathbf b$: shifted};
    \node[align=center] at (0,1.9) {\small same straightness,\\ \small but $\mathbf 0$ moved $\Rightarrow$ not linear};
  \end{scope}
  % nonlinear: bent
  \begin{scope}[xshift=10cm]
    \draw[->,gray] (-1.5,0)--(1.5,0); \draw[->,gray] (0,-1.2)--(0,1.4);
    \draw[amber,very thick] plot[domain=-1.1:1.1,smooth] (\x,{0.6*\x*\x});
    \node at (0,-1.7) {\small nonlinear: lines bend};
    \node[align=center] at (0,1.9) {\small fails even\\ \small before checking $\mathbf 0$};
  \end{scope}
\end{tikzpicture}

\caption{Linear vs. affine vs. nonlinear. Notice that affine maps are \emph{almost} linear — same flatness, wrong origin — which is exactly why they matter: a neural network layer is affine, and every ``bias'' term you will ever add is the $+\mathbf b$ that makes it non-linear in this strict sense.}
\label{fig:affine}
\end{figure}

\begin{example}\label{ex:linverify}
Is $T(x,y)=(x+2y,\;3x-y)$ linear? If so, find its matrix.
\textbf{Strategy.} Check the two rules symbolically, not on test values — and then read the matrix off what $T$ does to $\mathbf e_1,\mathbf e_2$: the columns \emph{are} the outputs.
\begin{solution}
$T(\mathbf u+\mathbf v)=(u_1+v_1+2(u_2+v_2),\,3(u_1+v_1)-(u_2+v_2))=T(\mathbf u)+T(\mathbf v)$ \checkmark; $T(c\mathbf v)=(cv_1+2cv_2,\,3cv_1-cv_2)=cT(\mathbf v)$ \checkmark. Linear. $T(\mathbf e_1)=(1,3)$, $T(\mathbf e_2)=(2,-1)$, so
\(
[T]=\begin{bmatrix}1&2\\3&-1\end{bmatrix}.
\)
\end{solution}
\end{example}

\section{Kernel and image}\label{sec:kerim}

Two subspaces accompany every linear map, and they are the Chapter-3 pair in disguise: the \emph{kernel}\index{kernel} $\ker T=\{\mathbf v:T(\mathbf v)=\mathbf 0\}\subseteq V$ (what the map erases) and the \emph{image}\index{image} $\operatorname{im}T=\{T(\mathbf v):\mathbf v\in V\}\subseteq W$ (what it can produce). For $T(\mathbf x)=A\mathbf x$ these are exactly $\ker(A)$ and $\operatorname{col}(A)$ — same objects, abstract names. Rank--nullity survives the translation with zero extra work:
\[
  \dim\operatorname{im}T + \dim\ker T = \dim V.
\]
Two names for the borderline cases: $T$ is \emph{injective}\index{injective} (no two inputs share an output) iff $\ker T=\{\mathbf 0\}$; \emph{surjective}\index{surjective} (reaches all of $W$) iff $\operatorname{im}T=W$. When both hold, $T$ is an \emph{isomorphism}\index{isomorphism} — the spaces are the same space with different labels, and $T^{-1}$ exists as a linear map.

\begin{figure}[htbp]
\centering
\begin{tikzpicture}[scale=1.0, >=stealth]
  \draw[fill=graybg!25,draw=gray] (0,0) ellipse (1.7 and 1.1);
  \node at (0,1.35) {$V$};
  \draw[fill=amberbg!50,draw=amber] (0,-0.15) ellipse (0.55 and 0.45);
  \node[amber] at (0,-0.15) {$\ker T$};
  \fill (0,-0.15) circle (1pt);
  \draw[fill=graybg!25,draw=gray] (6.5,0) ellipse (1.7 and 1.1);
  \node at (6.5,1.35) {$W$};
  \draw[fill=accentsoft!50,draw=accent] (6.5,-0.05) ellipse (0.8 and 0.6);
  \node[accent] at (6.5,-0.05) {$\operatorname{im} T$};
  \fill (6.5,-0.75) circle (1pt) node[below, ink]{\tiny $\mathbf 0_W$};
  \draw[->,thick] (1.9,0.3) -- (4.6,0.3) node[midway,above]{$T$};
  \draw[->,thick,dashed,amber] (0.9,-0.6) .. controls (2.5,-1.3) and (4,-1.3) .. (5.6,-0.85);
  \node[amber,align=center] at (3.2,-1.45) {\small $\ker T \mapsto \mathbf 0_W$};
\end{tikzpicture}

\caption{Kernel and image in one schematic. Notice the kernel is a subspace \emph{of the domain} that all lands on one point, while the image is a subspace \emph{of the codomain} that may be strictly smaller than $W$ — the two failures are independent, and rank--nullity counts them together.}
\label{fig:kerim}
\end{figure}

\section{Every linear map is a matrix}\label{sec:matrixrep}

Here is the theorem that justifies the whole book so far. Given a linear $T:V\to W$ and bases $\mathcal B=\{\mathbf b_1,\dots,\mathbf b_n\}$ of $V$, $\mathcal C=\{\mathbf c_1,\dots,\mathbf c_m\}$ of $W$: write each $T(\mathbf b_j)$ in $\mathcal C$-coordinates, and place those coordinates in column $j$. The resulting $m\times n$ matrix $[T]_{\mathcal{C}\leftarrow\mathcal{B}}$ \emph{is} $T$ in coordinates:
\[
  [T(\mathbf v)]_{\mathcal C} = [T]_{\mathcal{C}\leftarrow\mathcal{B}}\,[\mathbf v]_{\mathcal B}.
\]
The columns answer one question each — ``where does basis vector $\mathbf b_j$ go, in $\mathcal C$-words?'' — and that is the entire content of the map, because linearity extends basis behavior to everything. When both bases are standard, $[T]$ is just the matrix whose columns are $T(\mathbf e_j)$ — which is why Chapter 2 could find a matrix by watching $\mathbf e_j$.

Composition now earns its keep: $[S\circ T]=[S][T]$ — the matrix of a composed map is the product of the matrices, in whatever consistent bases you chose. \emph{That} is why matrix multiplication is defined the way it is (Figure~\ref{fig:composition}); Chapter 2's rule was reverse-engineered from this requirement.

\begin{figure}[htbp]
\centering
\begin{tikzpicture}[scale=1.0, >=stealth]
  \node[draw,graybg!40,fill=graybg!20,minimum width=1.5cm,minimum height=1cm] (U) at (0,0) {$U$};
  \node[draw,graybg!40,fill=graybg!20,minimum width=1.5cm,minimum height=1cm] (V) at (4,0) {$V$};
  \node[draw,graybg!40,fill=graybg!20,minimum width=1.5cm,minimum height=1cm] (W) at (8,0) {$W$};
  \draw[->,thick] (U) -- (V) node[midway,above]{$T$\quad$[T]$};
  \draw[->,thick] (V) -- (W) node[midway,above]{$S$\quad$[S]$};
  \draw[->,thick,accent] (U.south) .. controls (2,-1.3) and (6,-1.3) .. (W.south) node[midway,below]{$S\circ T$\quad$[S][T]$};
  \node[align=center] at (4,-2.1) {\small matrix multiplication is composition written in coordinates:\\ \small apply $[T]$'s columns through $[S]$};
\end{tikzpicture}

\caption{Composition is matrix multiplication. Notice the direction of travel: $\mathbf x$ meets $T$ (and $[T]$) first. The product $[S][T]$ is exactly the matrix whose columns are $[S]([T]\mathbf e_j)$ — do $T$, then $S$, to each basis vector.}
\label{fig:composition}
\end{figure}

\begin{example}\label{ex:gallery}
Find the matrices of the planar catalog: rotation by $\theta$, reflection across the $x$-axis, projection onto the $x$-axis.
\textbf{Strategy.} One method handles all three: watch $\mathbf e_1=(1,0)$ and $\mathbf e_2=(0,1)$, write their images as columns. Geometry supplies the images; trig supplies the numbers for rotation.
\begin{solution}
\emph{Rotation} $R_\theta$: $\mathbf e_1\mapsto(\cos\theta,\sin\theta)$, $\mathbf e_2\mapsto(-\sin\theta,\cos\theta)$ (derived properly in the hand derivation):
\(
R_\theta=\begin{bmatrix}\cos\theta&-\sin\theta\\ \sin\theta&\cos\theta\end{bmatrix}.
\)
\emph{Reflection} across the $x$-axis: $\mathbf e_1$ stays, $\mathbf e_2$ flips: $\begin{bmatrix}1&0\\0&-1\end{bmatrix}$.
\emph{Projection} onto the $x$-axis: $\mathbf e_1$ stays, $\mathbf e_2$ dies: $\begin{bmatrix}1&0\\0&0\end{bmatrix}$.
\end{solution}
\end{example}

\begin{example}\label{ex:diffpoly}
Differentiation on the space $P_3$ of cubics $a_0+a_1x+a_2x^2+a_3x^3$:
write $D=\frac{d}{dx}$ as a matrix.
\par\smallskip
\textbf{Strategy.} A map is linear if it respects $+$ and scaling ---
differentiation does both, so a matrix must exist. Feed it the basis:
$D(\mathbf{e}_j)$ is column $j$, exactly the
$A\mathbf{e}_j=\mathbf{a}_j$ trick of Chapter 2.
\begin{solution}
On the monomial basis $\mathcal{B}=\{1,x,x^2,x^3\}$:
$D(1)=0$, $D(x)=1$, $D(x^2)=2x$, $D(x^3)=3x^2$. Columns are those
images in $\mathcal{B}$-coordinates:
\[
  [D]_{\mathcal{B}}=
  \begin{bmatrix}0&1&0&0\\0&0&2&0\\0&0&0&3\\0&0&0&0\end{bmatrix}.
\]
Now $D(3x^3-2x)=(0,-2,0,9)^{T}$ by the matrix, and directly
$9x^2-2$ --- same. $[D]^4=0$: differentiating a cubic four times
annihilates it --- a nilpotent matrix doing honest work. \checkmark
\end{solution}
\end{example}

\section{Change of basis and similarity}\label{sec:changebasis}

Chapter 3 ended with coordinates as labels. Now the punchline: the \emph{same} linear map has different matrices in different bases, and the matrices are related by a rule with a name. If $\mathcal B$ is a basis of $V$ with basis-change matrix $B$ (the $\mathbf b_j$ as columns), then
\[
  [T]_{\mathcal B} = B^{-1}\,[T]_{\text{std}}\,B,
\]
because the right side does the whole trip: $B$ translates $\mathcal B$-coordinates to standard, $[T]_{\text{std}}$ does the map, $B^{-1}$ translates back. Two square matrices $M,N$ with $N=B^{-1}MB$ for some invertible $B$ are \emph{similar}\index{similar matrices} — and the meaning of similarity is exactly ``same map, different basis''. Chapter 8 exploits this: pick the basis where the map looks simplest (a diagonal matrix) and the map becomes transparent.

\begin{figure}[htbp]
\centering
\begin{tikzpicture}[scale=1.0, >=stealth]
  % standard-grid view
  \draw[gray!45,step=0.5] (-1.5,-1.5) grid (1.5,1.5);
  \draw[->] (0,0)--(1,0) node[right]{\small $\mathbf e_1$};
  \draw[->] (0,0)--(0,1) node[above]{\small $\mathbf e_2$};
  \draw[->,accent,very thick] (0,0)--(0,1) node[pos=0.5,left]{\tiny};
  \draw[->,accent,very thick] (0,0)--(-1,0);
  \node at (0,-1.9) {\small standard basis};
  \node[align=center] at (0,2) {\small quarter-turn: $\mathbf e_1\mapsto\mathbf e_2$,\\ \small $\mathbf e_2\mapsto-\mathbf e_1$};
  % B-grid view
  \begin{scope}[xshift=6.5cm]
    \foreach \s in {-1,-0.5,0,0.5,1}{\draw[gray!45] (\s*2+\s*1, \s*1+\s*1)+(-3,-2)--+(3,2);}
    \draw[->,accent,very thick] (0,0)--(2,1) node[right]{\small $\mathbf b_1$};
    \draw[->,accent,very thick] (0,0)--(1,1) node[above]{\small $\mathbf b_2$};
    \draw[->,good,very thick] (0,0)--(-1,2) node[above left=-2pt]{\small $T\mathbf b_1$};
    \draw[->,good,very thick] (0,0)--(-1,1) node[left]{\small $T\mathbf b_2$};
    \node at (0.5,-1.9) {\small basis $\mathcal B=\{(2,1),(1,1)\}$};
    \node[align=center] at (0.5,2) {\small same physical quarter-turn,\\ \small watched through $\mathcal B$-eyes};
  \end{scope}
\end{tikzpicture}

\caption{One quarter-turn, two witness bases. Notice the map itself did not move — only the grid did. The matrix $[T]_{\mathcal B}$ answers ``where do $\mathbf b_1,\mathbf b_2$ go, measured in $\mathbf b$'s?'' and similarity converts between the two descriptions.}
\label{fig:samemap}
\end{figure}

\begin{example}\label{ex:changebasis}
$T$ = quarter-turn of $\mathbb R^2$: $[T]_{\text{std}}=A=\begin{bmatrix}0&-1\\1&0\end{bmatrix}$. Find $[T]_{\mathcal B}$ for $\mathcal B=\{(2,1),(1,1)\}$, and verify it represents the same map.
\textbf{Strategy.} Apply the conversion: $[T]_{\mathcal B}=B^{-1}AB$ with $B=\begin{bmatrix}2&1\\1&1\end{bmatrix}$. Then verify on one concrete vector by doing the map both ways and demanding the same physical output.
\begin{solution}
$B^{-1}=\begin{bmatrix}1&-1\\-1&2\end{bmatrix}$, so
\(
[T]_{\mathcal B}=B^{-1}AB=\begin{bmatrix}-3&-2\\5&3\end{bmatrix}.
\)
Check on $\mathbf v=(5,1)$: standard gives $A\mathbf v=(-1,5)$. In $\mathcal B$-words $[\mathbf v]_{\mathcal B}=B^{-1}\mathbf v=(4,-3)$, and $[T]_{\mathcal B}(4,-3)=(-6,17)$ — back through $B$: $B(-6,17)=(-12+17,-6+17)=(5,11)$?? Careful — recompute: $B(-6,17)=(-6\cdot2+17\cdot1,\,-6+17)=(5,11)$. That disagrees! Recheck the arithmetic: $[T]_{\mathcal B}\,(4,-3)=\begin{bmatrix}-3&-2\\5&3\end{bmatrix}\begin{bmatrix}4\\-3\end{bmatrix}=\begin{bmatrix}-12+6\\20-9\end{bmatrix}=\begin{bmatrix}-6\\11\end{bmatrix}$, then $B(-6,11)=(-12+11,-6+11)=(-1,5)$ \checkmark — same as $A\mathbf v$. (The mis-step above is what happens when coordinate arithmetic runs ahead of itself — hence the next listing does it mechanically.)
\end{solution}
\end{example}

\begin{lstlisting}[language=Python,caption={Change of basis, done both ways on the same vector. File: \texttt{code/ch04\_change\_basis.py}.},label={lst:changebasis}]
import numpy as np
np.set_printoptions(precision=4, suppress=True)

# T = quarter-turn:  standard matrix A
A = np.array([[0., -1.], [1., 0.]])

# new basis B = {b1, b2} = columns of B
B = np.array([[2., 1.], [1., 1.]])      # b1=(2,1), b2=(1,1)

# same map in B-coordinates:  [T]_B = B^{-1} A B
TB = np.linalg.inv(B) @ A @ B
print("[T]_B =")
print(TB)

# check on a vector: v = (5, 1) standard.
v = np.array([5., 1.])
c = np.linalg.solve(B, v)             # [v]_B : B-coords of v
Tv_A = A @ v                          # [T(v)] in standard coords
Tv_B = B @ (TB @ c)                   # via B-coords
print("[v]_B =", c)
print("via A:", Tv_A)
print("via [T]_B:", Tv_B)
print("same physical answer:", np.allclose(Tv_A, Tv_B))
print("A = B [T]_B B^-1 ok:",
      np.allclose(A, B @ TB @ np.linalg.inv(B)))
\end{lstlisting}

\begin{lstlisting}[caption={Output of listing~\ref{lst:changebasis}. The two matrices look nothing alike yet produce identical physical outputs — similarity is sameness of map, not sameness of entries.}]
[T]_B =
[[-3. -2.]
 [ 5.  3.]]
[v]_B = [ 4. -3.]
via A: [-1.  5.]
via [T]_B: [-1.  5.]
same physical answer: True
A = B [T]_B B^-1 ok: True
\end{lstlisting}

\begin{exercise}[easy]
Find $[T]_{\text{std}}$ for $T(x,y,z)=(x+y,\;y-z,\;x+z)$ by watching $\mathbf e_1,\mathbf e_2,\mathbf e_3$.
\end{exercise}

\section{A geometric catalog: rotations, reflections, projections, shears}\label{sec:gallery}

Four planar maps deserve permanent residence (Figure~\ref{fig:gallery}); the code below moves a point cloud through each.

\begin{figure}[htbp]
\centering
\begin{tikzpicture}[scale=0.62, >=stealth]
\def\squaremarker{\draw[fill=accentsoft!50,draw=accent,thick] (0,0)--(1,0)--(1,1)--(0,1)--cycle; \node[accent] at (0.35,0.62){\tiny F};}
% original
\begin{scope}
  \draw[gray!45,step=0.5] (-0.5,-0.5) grid (1.5,1.5);
  \draw[->,gray] (-0.5,0)--(1.6,0); \draw[->,gray] (0,-0.5)--(0,1.6);
  \squaremarker
  \node at (0.5,-1) {\small original};
\end{scope}
% rotation 30deg
\begin{scope}[xshift=3.4cm]
  \draw[gray!45,step=0.5] (-0.5,-0.5) grid (1.5,1.5);
  \draw[->,gray] (-0.5,0)--(1.6,0); \draw[->,gray] (0,-0.5)--(0,1.6);
  \begin{scope}[rotate=30]\squaremarker\end{scope}
  \node at (0.5,-1) {\small rotation};
\end{scope}
% reflection y->-y
\begin{scope}[xshift=6.8cm]
  \draw[gray!45,step=0.5] (-0.5,-1.5) grid (1.5,0.5);
  \draw[->,gray] (-0.5,0)--(1.6,0); \draw[->,gray] (0,-1.6)--(0,0.6);
  \begin{scope}[yscale=-1]\squaremarker\end{scope}
  \node at (0.5,-1.9) {\small reflection};
\end{scope}
% shear x+=0.5y
\begin{scope}[xshift=10.2cm]
  \draw[gray!45,step=0.5] (-0.5,-0.5) grid (1.8,1.5);
  \draw[->,gray] (-0.5,0)--(1.9,0); \draw[->,gray] (0,-0.5)--(0,1.6);
  \begin{scope}[xslant=0.5]\squaremarker\end{scope}
  \node at (0.6,-1) {\small shear};
\end{scope}
% projection onto x-axis
\begin{scope}[xshift=13.8cm]
  \draw[gray!45,step=0.5] (-0.5,-0.5) grid (1.5,1.5);
  \draw[->,gray] (-0.5,0)--(1.6,0); \draw[->,gray] (0,-0.5)--(0,1.6);
  \draw[fill=amberbg!70,draw=amber,thick] (0,0)--(1,0)--(1,0.02)--(0,0.02)--cycle;
  \node[amber] at (0.45,0.2){\tiny F};
  \draw[dashed,gray] (1,1)--(1,0); \draw[dashed,gray] (0.35,0.62)--(0.35,0);
  \node at (0.5,-1) {\small projection};
\end{scope}
\end{tikzpicture}

\caption{The unit square (with a marker letter so you can see orientation) through the four classic maps. Notice: rotation preserves everything including handedness; reflection preserves distances but flips the letter; shear preserves area but not shape; projection collapses a whole dimension — its determinant of $0$ is the algebraic echo of the lost axis.}
\label{fig:gallery}
\end{figure}

\begin{description}
\item[Rotation] $R_\theta=\begin{bmatrix}\cos\theta&-\sin\theta\\ \sin\theta&\cos\theta\end{bmatrix}$: rigid, length- and angle-preserving; $\det=1$.
\item[Reflection] across a line: distance-preserving but orientation-reversing; $\det=-1$.
\item[Shear] $\begin{bmatrix}1&k\\0&1\end{bmatrix}$: slides each point horizontally by $k$ times its height — area-preserving, angle-destroying; $\det=1$.
\item[Projection] onto a line through the origin: $P=\mathbf{u}\mathbf{u}^{T}/\mathbf{u}^{T}\mathbf{u}$; idempotent ($P^2=P$), $\det=0$ — information genuinely lost.
\end{description}

\begin{lstlisting}[language=Python,caption={The four catalog maps applied to a point cloud, with determinants. File: \texttt{code/ch04\_gallery\_transform.py}.},label={lst:gallery}]
import numpy as np
np.set_printoptions(precision=3, suppress=True)

# the catalog maps as matrices
theta = np.pi / 6
rot = np.array([[np.cos(theta), -np.sin(theta)],
                [np.sin(theta),  np.cos(theta)]])
refl = np.array([[1., 0.], [0., -1.]])          # reflect across x-axis
shear = np.array([[1., 0.5], [0., 1.]])          # shear in x
proj = np.array([[1., 0.], [0., 0.]])            # project onto x-axis

# point cloud: the unit square + the letter F's bounding corners
pts = np.array([[0., 0.], [1., 0.], [1., 1.], [0., 1.],
                [0.25, 0.75], [0.75, 0.75]])

for name, M in [("rotation", rot), ("reflection", refl),
                ("shear", shear), ("projection", proj)]:
    print(f"{name:>12} det={np.linalg.det(M):+.3f}")
    print(np.round(pts @ M.T, 2))
print("rotation preserves lengths:",
      np.allclose(np.linalg.norm(pts, axis=1),
                  np.linalg.norm(pts @ rot.T, axis=1)))
print("projection is idempotent:", np.allclose(proj @ proj, proj))
\end{lstlisting}

\begin{lstlisting}[caption={Output of listing~\ref{lst:gallery} (first few rows shown per map). The determinants tell the story before the coordinates do: $+1$ preserves volume and orientation, $-1$ flips orientation, $0$ collapses.}]
    rotation det=+1.000
[[ 0.    0.  ]
 [ 0.87  0.5 ]
 [ 0.37  1.37]
 [-0.5   0.87]
 [-0.16  0.77]
 [ 0.27  1.02]]
  reflection det=-1.000
...
       shear det=+1.000
[[0.   0.  ]
 [1.   0.  ]
 [1.5  1.  ]
 [0.5  1.  ]
 [0.62 0.75]
 [1.12 0.75]]
  projection det=+0.000
[[0.   0.  ]
 [1.   0.  ]
 [1.   0.  ]
 [0.   0.  ]
 [0.25 0.  ]
 [0.75 0.  ]]
rotation preserves lengths: True
projection is idempotent: True
\end{lstlisting}

\begin{lstlisting}[language=Python,caption={Projection onto a general line, built from $P=\mathbf{u}\mathbf{u}^{T}/\mathbf{u}^{T}\mathbf{u}$ alone. File: \texttt{code/ch04\_projection\_matrix.py}.},label={lst:proj}]
import numpy as np
np.set_printoptions(precision=4, suppress=True)

# project onto the line spanned by u = (3, 4):  P = u u^T / u^T u
u = np.array([3., 4.])
P = np.outer(u, u) / (u @ u)
print("P =")
print(P)

x = np.array([7., 1.])
Px = P @ x
print("Px =", Px)
print("parallel to u?", np.allclose(Px, (Px @ u) / (u @ u) * u))
print("residual perpendicular?", np.allclose((x - Px) @ Px, 0))
print("P is symmetric:", np.allclose(P, P.T))
print("P^2 = P:", np.allclose(P @ P, P))

# every vector splits into along-u + perpendicular:
print("x = Px + (x-Px):", np.allclose(x, Px + (x - Px)))
\end{lstlisting}

\begin{lstlisting}[caption={Output of listing~\ref{lst:proj}. The projection of $(7,1)$ onto the line through $(3,4)$ lands exactly at $(3,4)$ — the nearest point on the line. Symmetric and idempotent are not coincidences; Chapter 7 shows they characterize orthogonal projections.}]
P =
[[0.36 0.48]
 [0.48 0.64]]
Px = [3. 4.]
parallel to u? True
residual perpendicular? True
P is symmetric: True
P^2 = P: True
x = Px + (x-Px): True
\end{lstlisting}

\begin{exercise}[medium]
$P=\mathbf{u}\mathbf{u}^{T}/\mathbf{u}^{T}\mathbf{u}$ for $\mathbf{u}=(1,2)$. Compute $P$, then $P(5,0)$ and $P(P(5,0))$ — and explain idempotence in one geometric sentence.
\end{exercise}

\begin{example}\label{ex:projmat}
Derive the matrix of orthogonal projection onto the line through
$\mathbf{u}=(1,2)^{T}$ --- then apply it to $\mathbf{v}=(3,0)^{T}$.
\par\smallskip
\textbf{Strategy.} Projection is linear (verify by picture: the
projected sum is the sum of projections), so find the matrix by feeding
it $\mathbf{e}_1,\mathbf{e}_2$ --- or shortcut via the outer-product
formula $P=\mathbf{u}\mathbf{u}^{T}/(\mathbf{u}\cdot\mathbf{u})$
(figure~\ref{fig:projection}).
\begin{solution}
$P=\frac{1}{5}\begin{bmatrix}1&2\\2&4\end{bmatrix}$. Check on
$\mathbf{v}$: $P\mathbf{v}=\tfrac15(3,6)^{T}=(0.6,1.2)^{T}$ --- a point
on the line through $(1,2)$ (it is $0.6\mathbf{u}$), and
$\mathbf{v}-P\mathbf{v}=(2.4,-1.2)^{T}$ is orthogonal:
$2.4\cdot1+(-1.2)\cdot2=0$. \checkmark\ Idempotence check:
$P^2=\frac1{25}\begin{bmatrix}5&10\\10&20\end{bmatrix}=P$ --- projecting
twice moves nothing, exactly as the geometry insists.
\end{solution}
\end{example}

\begin{figure}[htbp]\centering
% Figure 4.6: orthogonal projection onto a line through u.
\begin{tikzpicture}[scale=1.15,>=stealth]
  \draw[thick,gray] (-0.6,0) -- (3.4,1.7) node[right,font=\small] {line through $\mathbf{u}$};
  \draw[->,thick] (0,0) -- (0.95,0.5) node[midway,below,font=\small] {$\mathbf{u}$};
  \draw[->,thick] (0,0) -- (1.4,2.3) node[above,font=\small] {$\mathbf{v}$};
  \draw[->,thick,accent] (0,0) -- (1.71,0.86) node[below right=-2pt,font=\small] {$P\mathbf{v}=\frac{\mathbf u\cdot\mathbf v}{\mathbf u\cdot\mathbf u}\mathbf u$};
  \draw[dashed,gray] (1.4,2.3) -- (1.71,0.86);
  \draw (1.57,0.72) ++(0,0.12) -- ++(0.14,0.07) -- ++(-0.04,-0.14);
  \node[font=\footnotesize,align=center] at (3.4,2.4) {$P=\mathbf{u}\mathbf{u}^{T}$ when $\|\mathbf u\|=1$:\\ outer product \emph{is} the projector};
\end{tikzpicture}

\caption{Projection as a matrix. Notice the outer product doing the
work: $P\mathbf v$ keeps only the $\mathbf u$-component of $\mathbf v$.
The drop from $\mathbf v$ to $P\mathbf v$ is perpendicular --- always.}
\label{fig:projection}
\end{figure}

\section{Hand derivation}\label{sec:handderiv4}

Derive $R_\theta$ from trigonometry alone. Take $\mathbf e_1=(1,0)$: after rotation by angle $\theta$ it lands at $(\cos\theta,\sin\theta)$ — that is the definition of $\cos,\sin$ on the unit circle. Then $\mathbf e_2=(0,1)$ sits $90^\circ$ ahead, so it lands at $(\cos(\theta+90^\circ),\sin(\theta+90^\circ))=(-\sin\theta,\cos\theta)$. Assemble the columns. Then complete the check every derivation owes: verify $R_\theta R_\phi=R_{\theta+\phi}$ (the addition formulas emerge as a consequence), and $\det R_\theta=\cos^2\theta+\sin^2\theta=1$.

\section{Common mistakes}\label{sec:mistakes4}

\begin{tripwire}[Confusing an abstract vector with its coordinate vector]
$\mathbf v$ is the thing; $[\mathbf v]_{\mathcal B}$ is its description in $\mathcal B$-words. They coincide only in the standard basis, and writing ``$\mathbf v=(3,1)$'' without naming a basis is quietly assuming one. The bug appears the moment two bases are in play.
\end{tripwire}
\begin{tripwire}[Forgetting basis order matters]
$\{\mathbf b_1,\mathbf b_2\}$ and $\{\mathbf b_2,\mathbf b_1\}$ are different bases: coordinates and matrix columns come out swapped. A basis is an \emph{ordered} list; permute it and every coordinate vector permutes with it.
\end{tripwire}
\begin{tripwire}[{Applying $[T]_{\mathcal B}$ to coordinates in the wrong basis}]
$[T]_{\mathcal B}$ only eats $\mathcal B$-coordinate vectors. Feed it standard coordinates and it will cheerfully produce a wrong answer that still looks like a vector — the most dangerous kind. Before every matrix application, ask: which basis are these numbers in?
\end{tripwire}
\begin{tripwire}[Treating similar matrices as identical]
$A$ and $B^{-1}AB$ represent the same map but are different matrices: different entries, different pivots, different easy computations. Similarity preserves eigenvalues, determinant, rank — not entries. ``Same map'' is an equivalence, not an equality.
\end{tripwire}

\section{Summary}\label{sec:ch4summary}

\begin{center}
\small
\begin{tabularx}{\linewidth}{@{}lY Y@{}}
\toprule
\tablehead \thead{Concept} & \thead{Statement} & \thead{Watch for} \\
\midrule
Linear map & $T(a\mathbf u+b\mathbf v)=aT\mathbf u+bT\mathbf v$ & must fix $\mathbf 0$ \\
Kernel & $\{\mathbf v:T\mathbf v=\mathbf0\}\subseteq V$ & injective iff trivial \\
Image & $T(V)\subseteq W$ & surjective iff $=W$ \\
Rank--nullity & $\dim\operatorname{im}+\dim\ker=\dim V$ & ch.\,3 theorem, restated \\
$[T]_{\mathcal C\leftarrow\mathcal B}$ & columns $=[T\mathbf b_j]_{\mathcal C}$ & eats $\mathcal B$-coords only \\
Composition & $[S\circ T]=[S][T]$ & why multiplication exists \\
Change of basis & $[T]_{\mathcal B}=B^{-1}[T]_{\mathrm{std}}B$ & $B$ translates first \\
Similarity & same map, different basis & entries differ; invariants match \\
Catalog & rot ($\det1$), refl ($-1$), shear ($1$), proj ($0$, $P^2=P$) & projection onto $\mathbf u$: $\mathbf u\mathbf u^T/\mathbf u^T\mathbf u$ \\
\bottomrule
\end{tabularx}
\end{center}

\begin{wins}
You can certify a map linear, write its matrix in any basis, shuttle the same map between bases with similarity, and deploy the four canonical planar maps — the vocabulary every linear-model picture uses.
\end{wins}

\chapternotes{%
The definition-first register (linear map before matrix, kernel/image before rank--nullity) follows Axler chapter 3 \cite{axler}; our matrix-first order is Strang's, and this chapter is where the two conventions reconcile. The gallery treatment and grid-deformation figures borrow their pedagogy from \cite{ila} ch.9 and the \emph{Essence of Linear Algebra} series \cite{eola} (episodes 3--4: ``linear transformations'' and ``matrix multiplication as composition''). MML \S2.7 \cite{mml} covers the same ground in ML register.}
\begin{rulebox}[{Watch instead}]
\begin{enumerate}
\item \emph{Linear transformations and matrices} --- 3Blue1Brown,
Essence of Linear Algebra ch.~3, 10:58,
\url{https://youtu.be/kYB8IZa5AuE} --- the
grid-preserving view of \S\ref{sec:linmaps}--\S\ref{sec:matrixrep}
(accessed 2026-10-11).
\item \emph{Change of basis} --- 3Blue1Brown, Essence of Linear
Algebra ch.~13, 12:50, \url{https://youtu.be/P2LTAUO1TdA} ---
\S\ref{sec:changebasis} visualized (accessed 2026-10-11).
\end{enumerate}
\end{rulebox}

\section*{Further reading}
\addcontentsline{toc}{section}{Further reading}

\begin{itemize}
\item G.~Strang, \emph{Introduction to Linear Algebra}, 6th ed., \S7.1
(textbook level) --- linear transformations, matrix form, and
similarity at full length.
\item MIT OpenCourseWare, \emph{18.06SC Linear Algebra} --- OCW
recitations on matrix form of maps. \newline
OCW (course ``18.06SC''), at \url{https://ocw.mit.edu}
Verified 2026-10-11.
\item Margalit and Rabinoff, \emph{Interactive Linear Algebra},
``Linear Transformations'' and ``Matrix Transformations'' --- the same
machinery with more pictures. \newline
\url{https://textbooks.math.gatech.edu/ila/}
Verified 2026-10-11.
\item M.~Deisenroth et al., \emph{Mathematics for Machine Learning},
\S2.7 (free PDF) --- basis change in the ML register.\newline
\url{https://mml-book.github.io/book/mml-book.pdf}
Verified 2026-10-11.
\item J.~Hefferon, \emph{Linear Algebra}, ch.~III (free textbook) ---
maps defined rigorously, with the proof that matrix form is basis
dependence. \newline
\url{https://hefferon.net/linearalgebra/index.html}
Verified 2026-10-11.
\end{itemize}

\section*{Additional exercises}
\addcontentsline{toc}{section}{Additional exercises}

\begin{exercise}[easy]
Which of these maps $\mathbb{R}^2\to\mathbb{R}^2$ are linear? For each failure name the rule it breaks: (a) $(x,y)\mapsto(x+y,y)$; (b) $(x,y)\mapsto(xy,y)$; (c) $(x,y)\mapsto(x+1,y)$; (d) $(x,y)\mapsto(y,x)$.
\end{exercise}
\begin{exercise}[easy]
Find $[T]_{\text{std}}$: (a) $T$ = reflection across the line $y=x$; (b) $T$ = doubling every vector; (c) $T$ = projection onto the $y$-axis.
\end{exercise}
\begin{exercise}[easy]
$T:\mathbb{R}^3\to\mathbb{R}^3$ has matrix $\begin{bmatrix}1&0&1\\0&1&0\\0&0&0\end{bmatrix}$. Give a basis for $\ker T$ and $\operatorname{im}T$, and verify rank--nullity.
\end{exercise}
\begin{exercise}[easy]
Verify $[T][S]=[T\circ S]$ on $T(x,y)=(x+y,y)$, $S(x,y)=(2x,x-y)$: compute both sides, compare.
\end{exercise}
\begin{exercise}[medium]
Show $T:\mathcal{P}_2\to\mathcal{P}_2$, $T(p)=p'$ is linear, find its kernel and image, and write its matrix in the basis $\{1,t,t^2\}$. What is $[T]^3$ and why is it zero?
\end{exercise}
\begin{exercise}[medium]
$A=\begin{bmatrix}1&1\\0&1\end{bmatrix}$ (a shear). Compute $A^2,A^3,A^{10}$ and describe the map $T^k$ geometrically. What is $\lim_k A^k\mathbf e_1/\|A^k\mathbf e_1\|$?
\end{exercise}
\begin{exercise}[medium]
The quarter-turn $A=\begin{bmatrix}0&-1\\1&0\end{bmatrix}$ satisfies $A^2=-I$. Use only that fact (no trig) to compute $A^4, A^{2026}$, and the eigen-story it previews: can a real $2\times2$ matrix with $A^2=-I$ have a real eigenvector?
\end{exercise}
\begin{exercise}[medium — derive]
$\mathcal B=\{(1,1),(1,-1)\}$. For $T(x,y)=(x+y,x-y)$ find $[T]_{\mathcal B}$ two ways: (a) write $T\mathbf b_1,T\mathbf b_2$ in $\mathcal B$-coordinates directly; (b) compute $B^{-1}[T]_{\mathrm{std}}B$. Compare — they must agree.
\end{exercise}
\begin{exercise}[hard]
Prove $[S\circ T]_{\mathcal A\leftarrow\mathcal C}=[S]_{\mathcal A\leftarrow\mathcal B}\,[T]_{\mathcal B\leftarrow\mathcal C}$ for bases $\mathcal A,\mathcal B,\mathcal C$ — i.e.\ composition really is matrix multiplication in any consistent bases. (Expand both sides' $(i,j)$ entry; they are the same sum.)
\end{exercise}
\begin{exercise}[hard]
Prove similarity is an equivalence relation (reflexive, symmetric, transitive) and that similar matrices have the same rank and trace. Then exhibit $2\times2$ matrices with the same rank and trace that are \emph{not} similar — what extra invariant separates them?
\end{exercise}
\begin{exercise}[hard]
Implement change-of-basis as a function \texttt{to\_basis(A, B)} returning $B^{-1}AB$, then test it: generate random $2\times2$ $A$ and random bases $B$, verify $\operatorname{tr}$ and $\det$ are similarity-invariant across 100 trials. File: \texttt{code/ch04\_similarity\_invariants.py}.
\end{exercise}
\begin{exercise}[easy]
Is $T(p)=p'(0)$ linear on $P_3$? Write its matrix on the monomial basis
$\{1,x,x^2,x^3\}$. What are its kernel and image?
\end{exercise}
\begin{exercise}[medium]
For the projection matrix of example~\ref{ex:projmat}, compute
$P\mathbf{u}$ and explain in one sentence why the answer had to be
$\mathbf{u}$. Then project $\mathbf{v}=(1,1)^{T}$ and check
$\mathbf{v}-P\mathbf{v}\perp\mathbf{u}$.
\end{exercise}
\begin{exercise}[medium]
Write the matrix of the second derivative $D^2$ on $P_3$ in the
monomial basis, both by composing $[D]_{\mathcal B}$ with itself and by
differentiating each basis element twice. Confirm they match.
\end{exercise}
\begin{exercise}[hard]
Prove a projection matrix $P$ satisfies $P^2=P$ and $P^{T}=P$ iff it
is an orthogonal projection. Then show $I-P$ is also a projection, onto
the perpendicular line.
\end{exercise}
